\subsubsection{Examples}\label{subsec:distribution-derivative-examples}

\begin{Example}
\label{ex:heaviside-derivative}
Consider the function
\begin{equation}\label{eq:heaviside-function-definition}
H(x)=\begin{cases}0,&x\leq0\\1,&x>0.\end{cases}
\end{equation}
We compute
\begin{equation}\label{eq:heaviside-derivative-delta-pairing}
(H',\phi)=-(H,\phi')=-\int_0^\infty\phi'(x)\,\mathrm{d}x=\phi(0)=(\delta,\phi),
\end{equation}
i.e., the derivative of $H$ is the delta function. The function $H$ is called the \textbf{Heaviside function}.
\end{Example}

\begin{Example}
\label{ex:higher-delta-derivative}
The $k$th derivative of the delta function is the functional $\phi\mapsto(-1)^k\phi^{(k)}(0)$.
\end{Example}

\begin{Example}
\label{ex:singular-power-derivative}
Let
\begin{equation}\label{eq:positive-part-singular-power}
x_+^\lambda=\begin{cases}0,&x\leq0\\x^\lambda,&x>0\end{cases}
\end{equation}
and $-1<\lambda<0$. Naively, one may expect that the derivative is $\lambda x_+^{\lambda-1}$, but this function has a nonintegrable singularity and hence it is not a distribution. The proper answer is an appropriate regularization. We compute\[
\begin{aligned}
((x_+^\lambda)',\phi)&=-(x_+^\lambda,\phi')\\
&=-\int_0^\infty x^\lambda\phi'(x)\,\mathrm{d}x\\
&=-\lim_{\epsilon\to0}\int_\epsilon^\infty x^\lambda\phi'(x)\,\mathrm{d}x\\
&=\lim_{\epsilon\to0}\int_\epsilon^\infty\lambda x^{\lambda-1}(\phi(x)-\phi(\epsilon))\,\mathrm{d}x\\
&=\int_0^\infty\lambda x^{\lambda-1}(\phi(x)-\phi(0))\,\mathrm{d}x.
\end{aligned}
\]
\end{Example}

\begin{Example}
\label{ex:boundary-jump-derivative}
Let $\Omega$ be a domain with smooth boundary $\Gamma$. Let $f$ be in $C^1(\overline\Omega)$ and let $f=0$ in the exterior of $\Omega$. We regard $f$ as a distribution on $\mathbb{R}^m$. We find\[
\begin{aligned}
\left(\dfrac{\partial f}{\partial x_j},\phi\right)&=-\left(f,\dfrac{\partial\phi}{\partial x_j}\right)\\
&=-\int_\Omega f(\bm{x})\dfrac{\partial\phi}{\partial x_j}(\bm{x})\,\mathrm{d}\bm{x}\\
&=\int_\Omega\dfrac{\partial f}{\partial x_j}(\bm{x})\phi(\bm{x})\,\mathrm{d}\bm{x}
-\int_\Gamma f(\bm{x})\phi(\bm{x})n_j\,\mathrm{d}S.
\end{aligned}
\]Here $n_j$ is the $j$th component of the unit outward normal to $\Gamma$ and $\mathrm{d}S$ is differential $m-1$ dimensional surface area. Thus, the distributional derivative of $f$ involves one term corresponding to the ordinary derivative in $\Omega$ and another term involving a distribution supported on $\Gamma$. This latter term results from the jump of $f$ across $\Gamma$.
\end{Example}

\begin{Example}
\label{ex:polarization-charge}
The previous example has some important applications in electromagnetism. Let $\Omega\subset\mathbb{R}^3$ be a domain with smooth boundary $\Gamma$. Suppose we have a \textbf{polarization} vector field $\bm{p}:\overline\Omega\to\mathbb{R}^3$ which is in $C^1(\overline\Omega)$. By setting $\bm{p}=0$ in the exterior of $\Omega$, we can regard $\bm{p}$ as a distribution on $\mathbb{R}^3$. We then define the \textbf{polarization charge} to be the divergence of $\bm{p}$ \emph{in the sense of distributions}. We calculate this as follows.\[
\begin{aligned}
(\nabla\cdot\bm{p},\phi)&=-\sum_{i=1}^3\left(p_i,\dfrac{\partial\phi}{\partial x_i}\right)\\
&=-\int_\Omega\bm{p}(\bm{x})\cdot\nabla\phi(\bm{x})\,\mathrm{d}\bm{x}\\
&=\int_\Omega\nabla\cdot\bm{p}(\bm{x})\phi(\bm{x})\,\mathrm{d}\bm{x}
-\int_\Gamma\bm{p}(\bm{x})\cdot\bm{n}(\bm{x})\phi(\bm{x})\,\mathrm{d}A.
\end{aligned}
\]Here $\bm{n}$ is the unit outward normal to $\Gamma$ and $\mathrm{d}A$ is differential surface area on $\Gamma$. Thus, the polarization charge involves one term corresponding to the ordinary divergence of $\bm{p}$ in $\Omega$ and surface charge given by the normal component of $\bm{p}$ on $\Gamma$. This latter term results from the jump of $\bm{p}$ across $\Gamma$. If $\bm{p}$ was piecewise smooth with surfaces of jump discontinuity in the interior of $\Omega$, the normal components of the jumps along these surfaces would contribute polarization charge as well.
\end{Example}

\begin{Example}
\label{ex:inverse-radius-laplacian}
Let $f(\bm{x})=1/|\bm{x}|=1/r$ on $\mathbb{R}^3$. It is easy to check that $\Delta(1/r)=0$ for $r\ne0$. We shall evaluate the Laplacian of $1/r$ in the distributional sense. We compute\[
\begin{aligned}
\left(\Delta\dfrac1r,\phi\right)&=\left(\dfrac1r,\Delta\phi\right)\\
&=\int_{\mathbb{R}^3}\dfrac{\Delta\phi(\bm{x})}r\,\mathrm{d}\bm{x}\\
&=\lim_{\epsilon\to0}\int_{r\geq\epsilon}\dfrac{\Delta\phi}r\,\mathrm{d}\bm{x}.
\end{aligned}
\]Integration by parts yields
\begin{equation}\label{eq:inverse-radius-laplacian-integration-by-parts}
\begin{aligned}
\int_{r\geq\epsilon}\dfrac{\Delta\phi}r\,\mathrm{d}\bm{x}
&=\int_{r\geq\epsilon}\Delta\left(\dfrac1r\right)\phi\,\mathrm{d}\bm{x}
-\int_{r=\epsilon}\dfrac{\partial\phi}{\partial r}\dfrac1r\,\mathrm{d}S\\
&\quad+\int_{r=\epsilon}\phi\dfrac{\partial}{\partial r}\dfrac1r\,\mathrm{d}S.
\end{aligned}
\end{equation}
On the right-hand side of \eqref{eq:inverse-radius-laplacian-integration-by-parts}, the first term vanishes, the second is of order $\epsilon$ as $\epsilon\to0$ and the last term is equal to $-\epsilon^{-2}\int_{r=\epsilon}\phi\,\mathrm{d}S$, i.e., to $-4\pi$ times the average of $\phi$ on the sphere of radius $\epsilon$. Letting $\epsilon\to0$, we therefore obtain
\begin{equation}\label{eq:inverse-radius-laplacian-delta-pairing}
\left(\Delta\dfrac1r,\phi\right)=-4\pi\phi(0),
\end{equation}
i.e.,
\begin{equation}\label{eq:inverse-radius-laplacian-delta}
\Delta\dfrac1r=-4\pi\delta.
\end{equation}
Solutions of the equation $Lu=\delta$, where $L$ is a partial differential operator with constant coefficients, are of considerable importance; we shall investigate more such solutions in the next two sections.
\end{Example}
